\chapter{Kidney Pain}

\section*{Definition and Overview}
Kidney pain is a discomfort felt in the flank, the area on either side of the spine between the lower ribs and the hip, or in the back just below the ribs. It is caused by conditions that stretch, inflame, or block the kidneys, most commonly kidney stones, kidney infection, and, less often, tumours or bleeding. Kidney stone pain is typically sudden, severe, and colicky, while infection pain is more constant with fever. Pain that is actually in the muscles, spine, or bowel can be mistaken for kidney pain, so assessment is important. Because some causes can damage the kidney or become life-threatening, persistent or severe flank pain should be assessed promptly.

\section*{Epidemiology}
Kidney stones affect around one in ten Australians at some time and are the most common cause of acute kidney pain, particularly in men and in people aged 30 to 60. Kidney infections are common in women and cause flank pain with fever. Kidney cancer, which can cause a dull flank ache, is less common and often silent. Back and muscle pain is far more common than true kidney pain and is the usual cause of flank discomfort in the community. Painless kidney disease is the rule, so kidney pain is not a reliable sign of chronic kidney disease.

\section*{Aetiology and Pathophysiology}
The kidney itself has no pain fibres in its filtering tissue; pain arises when the kidney capsule is stretched or the ureter is irritated. A kidney stone stuck in the ureter causes the ureter to spasm, producing waves of severe colicky pain that spread to the groin. Infection inflames the kidney, distending the capsule and causing constant flank pain with fever. Obstruction from any cause, including stones, tumours, or an enlarged prostate, stretches the kidney and causes a dull ache. Tumours and bleeding can stretch the capsule. The kidneys sit high in the back, so kidney pain is felt in the flank rather than the lower back, and it does not move with position or activity like muscle pain.

\section*{Clinical Features}
\begin{itemize}
    \item Pain in the flank, side, or upper back, on one or both sides
    \item Stone pain: sudden, severe, colicky waves radiating to the groin and genitals
    \item Infection pain: constant ache with fever, chills, and urinary symptoms
    \item Pain that is not changed by movement or position
    \item Blood in the urine
    \item Nausea and vomiting with severe pain
    \item A feeling of pressure or fullness in the flank
    \item Urgency, frequency, or burning with urination
    \item Pain from muscles and spine is more common and moves with movement
\end{itemize}

\section*{Differential Diagnosis}
The cause of flank pain must be identified.
\begin{itemize}
    \item \textbf{Kidney stones:} the most common cause of severe colicky pain
    \item \textbf{Kidney infection:} fever with constant flank pain
    \item \textbf{Kidney tumour or bleeding:} dull ache with blood in the urine
    \item \textbf{Muscle strain and back problems:} movement-related pain
    \item \textbf{Bowel conditions:} appendicitis, diverticulitis, and constipation
    \item \textbf{Gallbladder disease:} right upper abdominal pain
    \item \textbf{Shingles:} burning pain before a rash appears
    \item \textbf{Referred pain:} from the spine, hip, or abdominal organs
\end{itemize}

\section*{When to Seek Medical Care}
Sudden, severe flank pain, especially with blood in the urine or fever, warrants urgent assessment, and pain with vomiting, difficulty passing urine, or signs of sepsis requires emergency care. Persistent or recurrent flank pain should be assessed, and people with risk factors for stones or kidney disease should be tested. Do not assume flank pain is muscular without assessment.

\textbf{Call 000 (or local emergency services) for:}
\begin{itemize}
    \item Severe flank pain with fever, shaking, or confusion, which may indicate sepsis
    \item Inability to pass urine
    \item Vomiting with an inability to keep fluids down
    \item Blood in the urine with pain
    \item Collapse, dizziness, or severe illness
\end{itemize}

\section*{Diagnosis and Investigations}
\begin{itemize}
    \item \textbf{History and examination:} the site, nature, and radiation of the pain
    \item \textbf{Urine tests:} blood, protein, infection, and crystals
    \item \textbf{Blood tests:} kidney function, inflammatory markers, and electrolytes
    \item \textbf{Imaging:} CT scan is the best test for stones, with ultrasound for other causes
    \item \textbf{Pregnancy test:} in women of reproductive age
    \item \textbf{Assessment of the cause:} stones, infection, and structural problems
    \item \textbf{Follow-up:} confirming resolution
\end{itemize}

\section*{Management}
\begin{itemize}
    \item \textbf{Pain relief:} strong analgesics for stone pain
    \item \textbf{Fluids:} hydration to help pass small stones
    \item \textbf{Stone treatment:} observation for small stones, or procedures to break up or remove larger stones
    \item \textbf{Antibiotics:} for kidney infection
    \item \textbf{Treatment of the cause:} tumour, obstruction, and bleeding
    \item \textbf{Relief of obstruction:} stents or drainage where needed
    \item \textbf{Musculoskeletal care:} physiotherapy and pain management where the pain is muscular
    \item \textbf{Follow-up:} repeat imaging and tests
\end{itemize}

\section*{Complications}
\begin{itemize}
    \item Kidney damage from prolonged obstruction by a stone
    \item Sepsis from infected obstruction
    \item Abscess and kidney scarring from infection
    \item Recurrent stones
    \item Chronic pain and repeated presentations
    \item Missed serious conditions when pain is assumed to be muscular
\end{itemize}

\section*{Prognosis and Outcomes}
The outlook depends on the cause. Most kidney stones pass on their own or are removed successfully, with full recovery, and kidney infections resolve with antibiotics. Tumours are best treated when found early. Muscular pain settles with physiotherapy and activity modification. The key is accurate diagnosis, because treating the right cause resolves the pain and prevents complications.

\section*{Prevention and Self-Care}
\begin{itemize}
    \item Drink plenty of fluids, which is the best prevention for stones
    \item Limit salt and reduce foods that promote stones if you are prone
    \item Treat urinary infections promptly
    \item Seek early assessment for flank pain
    \item Attend follow-up imaging after stones
    \item Manage weight and stay active
    \item Report blood in the urine promptly
    \item Seek urgent care for fever with pain
\end{itemize}

\section*{Sources and Further Reading}
The following reputable sources provide further information for healthcare students and patients seeking reliable, current advice about kidney pain.
\begin{itemize}
    \item Kidney Health Australia. \textit{Kidney pain and stones.} https://kidney.org.au/
    \item Healthdirect Australia. \textit{Kidney pain.} https://www.healthdirect.gov.au/
    \item Better Health Channel. \textit{Kidney stones.} https://www.betterhealth.vic.gov.au/
    \item Urological Society of Australia and New Zealand. \textit{Kidney stone information.} https://www.usanz.org.au/
    \item NHS. \textit{Kidney stones and kidney pain.} https://www.nhs.uk/
    \item Mayo Clinic. \textit{Kidney pain.} https://www.mayoclinic.org/
\end{itemize}
